\section{Caso de fracaso: por qué la primera arquitectura con Kubernetes no era adecuada}

La primera versión de Vercel usó Kubernetes como plataforma de ejecución de cargas de trabajo. Kubernetes en sí no era el error; el problema era que la promesa del producto consistía en que «cada Git commit se convierta en un despliegue rápido, inmutable y compartible». Si cada deployment se asigna directamente a un pod o a una unidad de recursos estáticos, el número de despliegues crece con los commits, y la idle reservation, la latencia de planificación y las scale transitions destruyen la economía unitaria.

\subsection{Que algo sea técnicamente viable no significa que sea económicamente viable}

Supongamos que en un día se generan $N_d$ preview deployments, que el costo de los recursos ociosos de cada deployment es $c_{idle}$ y que el costo de ejecución real es $c_{run}$; entonces el costo diario aproximado es
\[
C_{day}=N_d\cdot c_{idle}+\sum_j c_{run,j}.
\]
Cuando la mayoría de las previews casi no recibe visitas, el primer término domina el costo. El sistema debe lograr que los «despliegues sin visitas» tengan un costo de ejecución cercano a cero y, al mismo tiempo, se reactiven rápidamente cuando llega la primera solicitud; solo así puede sostener el per-commit product model.

\lecturefigure{02-kubernetes-economics.png}{La promesa de producto per-commit amplifica el costo y la latencia de la asignación estática de recursos.}{Subtítulos locales 01:12--01:42, 28:50--30:10; redibujo conceptual.}

\begin{knowledgebox}{Lectura de la figura: el problema no es el pod, sino la correspondencia}
Si cada immutable deployment queda ligado de forma permanente a un pod independiente, el número de previews y los recursos en ejecución crecen de manera aproximadamente lineal. Una arquitectura más adecuada representa el despliegue como artifact + metadata y, cuando llega una solicitud, reutiliza el runtime, la cache y los recursos de planificación compartidos. Kubernetes puede servir como componente de bajo nivel, pero los objetos del producto no deben quedar ligados uno a uno con los objetos de recursos.
\end{knowledgebox}

\begin{warningbox}{«Kubernetes es demasiado pesado» no es una conclusión general}
La valoración de la clase se refiere a la cantidad altísima de despliegues que Vercel tenía entonces, al scale-to-zero rápido y al modelo económico de las previews. Para servicios de ejecución prolongada, capacidad estable y equipos de plataforma maduros, Kubernetes todavía puede ser una opción razonable. El diseño de sistemas debe comparar la workload shape, la escala del plano de control y el costo unitario, en lugar de repetir preferencias por una herramienta.
\end{warningbox}

\teachervoice{Rauch describe este fracaso sin rodeos: la primera implementación estaba «bleeding money left and right». El valor de esta teacher voice está en que no rebajó el estándar de producto de «desplegar cada commit», sino que convirtió una meta de experiencia poco razonable en nuevos problemas de investigación sobre scheduler, placement y metadata.}

\subsection{Resumen de la sección}

La arquitectura debe ajustarse a la cardinalidad y al ciclo de vida de los objetos del producto. El per-commit deployment requiere artifact-first, shared-runtime y fast activation, en lugar de asignar de forma permanente cada despliegue lógico a un arrendamiento de recursos estáticos.
